\section{浮置导体的消元与机械功}
\subsection{给定电量的导体不能直接当作接地端}
多导体系统可写$\vct Q=\mathsf C\vct V$。以足够明确的参考导体固定零势后，把导体分成电势给定组$a$和电量给定的浮置组$f$，有
\[
 Q_f=C_{fa}V_a+C_{ff}V_f.
\]
若浮置子块可逆，先解$V_f=C_{ff}^{-1}(Q_f-C_{fa}V_a)$，再代回受控端：
\[
 Q_a=(C_{aa}-C_{af}C_{ff}^{-1}C_{fa})V_a
 +C_{af}C_{ff}^{-1}Q_f.
\]
第一项是消去浮置自由度后的电容矩阵，第二项保留它的预存电量。将浮置导体直接置为$V_f=0$，等于额外接入电源，物理问题已经改变。中性浮置导体只满足总电量零，两个表面仍可带相反电荷。

求机械力还要分别处理控制方式。全体电量固定时，取场能$U=\tfrac12\vct V^{\mathsf T}\mathsf C\vct V$并对几何坐标微分；某些电势由理想电源保持时，适合使用$\mathcal G=U-\sum_{a}V_aQ_a$，在固定受控电势及浮置电量下取$F=-\partial\mathcal G/\partial x$。减去的项记录电源交换的能量。

\begin{example}[带电浮置板的电势与失稳]
两块面积$A$的大平行板位于$x=0,d$，电势分别为0、$V$。中间放一块极薄导体板，位置$x$，总电量固定为$Q$，不接电源。间隙为真空，忽略边缘效应。求中间板电势、电力及其可能的平衡位置，并分别用压力和能量法求力。
\ExpansionDiagram{fig:dep-es3-floating-q}{两侧固定电势，中间板固定总电量}{
\draw[thick](0,-1)--(0,1.5)node[above]{$0$};\draw[very thick](2,-1)--(2,1.5)node[above]{$v,Q$};\draw[thick](5,-1)--(5,1.5)node[above]{$V$};\draw[->](-.3,-1.4)--(5.7,-1.4)node[right]{$x$};\draw[<->](0,-.5)--(2,-.5)node[midway,below]{$x$};\draw[<->](2,-.5)--(5,-.5)node[midway,below]{$d-x$};}
\end{example}
\begin{solution}
令$K=\varepsilon_0A$。图\ref{fig:dep-es3-floating-q}中左右间隙分别等效为$C_1=K/x$、$C_2=K/(d-x)$。中间板两面的电量之和为
\[
 Q=C_1v+C_2(v-V),\qquad
 v=\frac{Q+C_2V}{C_1+C_2}
 =\frac{Vx}{d}+\frac{Qx(d-x)}{Kd}.
\]
沿$+x$定义有向电场，左右两段分别为
\[
 E_L=-\frac vx=-\frac Vd-\frac{Q(d-x)}{Kd},\qquad
 E_R=-\frac{V-v}{d-x}=-\frac Vd+\frac{Qx}{Kd}.
\]
压力作用在板的两个外法线方向，故净力向右取正：
\begin{align*}
 F&=\frac K2(E_R^2-E_L^2)\\
 &=-\frac{QV}{d}+\frac{Q^2(2x-d)}{2Kd}.
\end{align*}
\ExpansionDiagram{fig:dep-es3-floating-a}{两侧电压力相减形成净力；平衡附近的位移会被放大}{
\draw[very thick](2,-1)--(2,1.2);\draw[->,thick](2,.5)--(.5,.5)node[left]{$KE_L^2/2$};\draw[->,thick](2,-.3)--(3.5,-.3)node[right]{$KE_R^2/2$};\draw[->](4.8,-1)--(7,-1)node[right]{$x-x_e$};\draw[->](5.8,-2.2)--(5.8,1.2)node[above]{$F$};\draw[thick](4.9,-1.9)--(6.8,0);}
图\ref{fig:dep-es3-floating-a}中的力随$x$线性增加。若$Q\ne0$，平衡点为$x_e=d/2+KV/Q$，只有$0<x_e<d$时才在间隙内。由于$\partial F/\partial x=Q^2/(Kd)>0$，该点不稳定；若$Q=0$，则各处净力均为零，属于忽略边缘效应后的中性平衡。

能量法要保留右侧电源的功。右极板电量为$Q_R=C_2(V-v)=KV/d-Qx/d$。两段场能为
\[
 U=\frac12C_1v^2+\frac12C_2(V-v)^2
 =\frac{KV^2}{2d}+\frac{Q^2x(d-x)}{2Kd}.
\]
只微分$U$会遗漏$-QV/d$一项。正确受控势函数为
\[
 \mathcal G=U-VQ_R=-\frac{KV^2}{2d}
 +\frac{QVx}{d}+\frac{Q^2x(d-x)}{2Kd}.
\]
在固定$Q,V$下，$\partial_x[x(d-x)]=d-2x$，所以
\[-\partial_x\mathcal G=-\frac{QV}{d}-\frac{Q^2(d-2x)}{2Kd},\]
恰为压力法的净力。
\end{solution}

\subsection{浮置消元怎样保留串联电容的物理含义}
若上题中间板中性，$Q=0$，有$v=C_2V/(C_1+C_2)$，外部等效电容$C_1C_2/(C_1+C_2)=K/d$。中间极薄板改变两段电压分配，却不改变总容量。若板厚为$t$，有效真空间隙总长为$d-t$，容量改为$K/(d-t)$；若中间板接地，右侧电源只看到$K/(d-x)$。这三种模型分别对应中性浮置、有限厚度浮置和接地，不能仅凭几何图上“多了一块导体”判断容量增加多少。

预存电量$Q\ne0$时，右极板电量对$V$的斜率仍是$K/d$，电荷—电压直线却多出截距$-Qx/d$。容量只给出增量响应；初始电场、机械力和储能还取决于预存电量。多导体系统中的浮置金属件也应保留各自的电量条件。
